\documentclass{article}
\usepackage{amsmath}
\usepackage{amsthm}
\usepackage{amssymb}

\newtheorem{theorem}{Theorem}
\newtheorem{lemma}[theorem]{Lemma}
\newtheorem{proposition}[theorem]{Proposition}
\newtheorem{corollary}{Corollary}[theorem]
\newtheorem{definition}{Definition}
\newtheorem{remark}{Remark}

\begin{document}

\section{The Fractal Tennis Computing Architecture (FTCA)}

To realize the emergent quantum mechanics of FSEP in a computational substrate, we introduce the \emph{Fractal Tennis Computing Architecture (FTCA)}, a geometric self-referential hyperparallel computing system based on the same Apollonian sphere-packing and M\"obius-group algebra. The FTCA embodies the persistent fractal boundaries of FSEP in hardware, enabling efficient classical simulation of quantum circuits through distributed, scale-recursive processing.

\begin{remark}[Layer Separation]
The results in this section are organized in three distinct layers: (1)~a classical simulation theorem establishing that FTCA efficiently simulates arbitrary quantum circuits; (2)~an emergent-QM corollary tying the internal optical dynamics to the FSEP ontology; and (3)~the MATHICCS/FSEP background, which motivates the architecture but is not required for the simulation theorem itself.
\end{remark}

\subsection{Geometric Structure and Sphere Taxonomy}

The FTCA consists of tangent polytope-spheres arranged in an Apollonian-like packing, starting from a tetrahedral seed: an inner core sphere and four mutually tangent larger spheres, replicated iteratively outward. Spheres are classified as:
\begin{itemize}
\item \textbf{T2 Balls}: Leaf spheres with radii from a finite basis set $\mathcal{R}_{\text{basis}} = \{r_1, \dots, r_N\}$ (determined by the seed). These host \emph{Heterogeneous Node Tiles (HNTs)} on polygonal faces, each HNT comprising FPGA/CPU/GPU substrates, a $W \times H$ 2D optical emitter/sensor array, and an analog op-amp comparator bank.
\item \textbf{Virtual Outer Container (VOC) Balls}: Composite spheres with radii $\notin \mathcal{R}_{\text{basis}}$, containing sub-assemblies of T2/VOC Balls. The outermost is the \emph{Virtual Outer Soddy (VOS)}, acting as a higher-order seed for fractal extension.
\end{itemize}

Each sphere's identity is encoded by an \emph{Apollonian Group Address (AGA)} $(k, w)$, where $k$ is the Descartes curvature (encoding radius $r = 1/|k|$) and $w$ is a prime-product word from generators $S_1 \to 2$, $S_2 \to 3$, $S_3 \to 5$, $S_4 \to 7$. The \emph{depth} of a sphere is the number of prime factors in $w$ counted with multiplicity. At fractal depth $n$, there are $\Theta(k^n)$ T2 Balls, where $k \approx 4$--$6$ is the Apollonian branching factor.

Interconnects include tangency links (electrical/optical/fiber tiers) and \emph{tubes} (void-routed fiber bundles via Tangent Communication Modules). A \emph{Connectivity Discovery Record (CDR)} tracks links, including \emph{phantom entries} for predicted but uninstalled connections, with geometric delays $\tau_{\text{geom}} = \sum 2r_i / c_{\text{med}}$.

\subsection{Atomic Fractal Logic Gate (AFLG)}

The core primitive is the \emph{Atomic Fractal Logic Gate (AFLG)}, executed per pixel on HNT optical arrays:
\[
\text{Sense} \to \text{Evaluate}(f(\mathbf{t})) \to \text{Act} \to \text{Wait}(\tau) \to \text{Recalculate},
\]
where $\mathbf{t} = (\lambda, \theta, \phi, I)$ is a trigger vector from op-amp comparators. AFLGs replicate fractally across scales (pixel $\to$ tile $\to$ T2 Ball $\to$ VOC Ball $\to$ VOS), mirroring FSEP's recursive boundaries.

HNTs expose a \emph{Tile Interface Contract (TIC)} for capability advertisement (including AGA) and workload dispatch. Channels are managed via a \emph{Channel Assignment Table (CAT)}, with three-tier medium promotion for adaptive routing.

\subsection{Formal Computational Model}

We now define FTCA as a precise computational model so that the simulation theorem can be stated and proved.

\begin{definition}[FTCA Configuration and Global Step]
Fix a branching factor $k$ and fractal depth $n$. A \emph{configuration} $\mathcal{S}_n$ of FTCA at depth $n$ is the tuple of all tile registers, op-amp DAC thresholds, and CAT/CDR states on every T2 Ball at depth $\le n$. A \emph{global step} is one synchronous AFLG cycle (sense--evaluate--act--wait--recalculate) executed in parallel on every pixel of every HNT at all depths $\le n$, together with the induced updates to CDR/CAT on ADMIN and QUANTUM\_STATE channels.

A \emph{program} on FTCA consists of:
\begin{enumerate}
  \item An initial configuration $\mathcal{S}_n^{(0)}$ (encoding the input state $|\psi_0\rangle$).
  \item A finite sequence of AFLG rule tables and DAC threshold patterns $\{\mathcal{R}_\ell\}_{\ell=1}^L$ applied layerwise, loaded via AFLGRULELOAD CAT channels before each global step.
\end{enumerate}
The output is read from the QUANTUM\_STATE channel after $L$ global steps.
\end{definition}

\subsection{Hilbert-Space Encoding}

\begin{lemma}[Classical State Embedding]
\label{lem:embedding}
Let $w$ be the number of logical qubits. Choose fractal depth $n = \Theta(w / \log_2 k)$, so that the number of T2 Balls at depth $\le n$ satisfies $\Theta(k^n) \ge 2^w$. There exists an injective map
\[
\iota: \{0, \dots, 2^w - 1\} \to \{\text{T2 Balls at depth} \le n\}
\]
assigning each computational basis state $|i\rangle$ to one distinct T2 Ball (equivalently, to a tile register on that ball). Using the HNT's local registers and AFLG internal state, define a memory layout such that the quantum state
\[
|\psi\rangle = \sum_{i=0}^{2^w-1} \alpha_i |i\rangle
\]
is represented by storing $\alpha_i \in \mathbb{C}$ (in fixed- or floating-point) in tile $\iota(i)$'s registers, exposed on the QUANTUM\_STATE channel for global normalization via echo-shell aggregation.
\end{lemma}

\begin{proof}[Proof sketch]
Injectivity follows from the capacity bound: there are $\Theta(k^n) \ge 2^w$ T2 Balls at depth $\le n$, so any injective assignment of basis indices to balls exists. One concrete choice interprets the lower $w$ bits of the prime-multiset word of each AGA as the binary expansion of the basis index $i$. Global normalization is performed via the existing echo-aggregation protocol on QUANTUM\_STATE CAT channels, which sums $|\alpha_i|^2$ across all tiles and normalizes in $O(\log(k^n)) = O(n)$ rounds.
\end{proof}

\subsection{Gate Implementation via AFLG 5-Step Cycle}

\begin{lemma}[Gate Gadgets]
\label{lem:gates}
Let $\mathcal{G} = \{H, T, \mathrm{CNOT}\}$ be the standard universal gate set. For any $\varepsilon > 0$, there exist a finite AFLG rule table and DAC threshold configuration such that:
\begin{enumerate}
  \item \textbf{Single-qubit gates}: For any tile carrying amplitude pair $(\alpha, \beta)$, one global AFLG cycle applies an analog rotation---implemented via DAC-updated thresholds on polarization/intensity comparators, with gate parameters broadcast as LED standing-wave patterns from the admin core---that transforms $(\alpha, \beta)$ to within $\varepsilon$ of $H(\alpha, \beta)$, $T(\alpha, \beta)$, or identity.
  \item \textbf{Two-qubit gates}: For any pair of tiles $A$ (control) and $B$ (target) connected by a tangency link or tube, a constant number of AFLG cycles plus a single tangency message implements the CNOT update on the four-amplitude tuple $(\alpha_{00}, \alpha_{01}, \alpha_{10}, \alpha_{11})$, again up to error $\varepsilon$.
  \item \textbf{Parallel execution}: An HNT with $W \times H$ pixels and $E$ edges supports up to $WHE$ independent AFLGs per clock; hence an entire gate layer acting on all $2^w$ amplitude tiles executes in a single global step when $n \ge w / \log_2 k$.
\end{enumerate}
\end{lemma}

\begin{proof}[Proof sketch]
Single-qubit gates: An arbitrary $U \in \mathrm{SU}(2)$ is $\varepsilon$-approximated by a sequence of $H$ and $T$ gates (Solovay--Kitaev theorem), each of which is implemented by a suitable DAC threshold shift and optical broadcast on the AFLGRULELOAD channel. The analog error per gate is bounded by optical SNR and DAC resolution, giving $\varepsilon_{\text{gate}} \le 2^{-p}$ for $p$-bit precision.

Two-qubit CNOT: The control tile sends a conditioned optical burst (wavelength $\lambda'$ encodes the Pauli channel) over the tangency or tube link to the target tile. The target pixel evaluates the Boolean function $f(\mathbf{t}) = (\lambda', \theta, \phi, I)$ and updates $(\alpha_{10}, \alpha_{11}) \leftarrow (\alpha_{11}, \alpha_{10})$ via local FPGA logic. Echo return confirms completion; T3 fiber and CDR-based delay/timeout handling provide deterministic timing, absorbing all routing latencies into the Wait$(\tau)$ substep.

Parallelism: AFLG instances for non-overlapping qubit pairs act on disjoint tiles and execute simultaneously within one global step, exactly as in a parallel gate layer.
\end{proof}

\subsection{Universality and Resource Scaling}

\begin{proposition}[Resource Scaling]
\label{prop:scaling}
For any circuit on $w$ qubits and depth $d$, there exists a fractal depth $n = \Theta(w / \log_2 k)$ and an FTCA program at depth $n$ such that the circuit is simulated with:
\begin{itemize}
  \item \textbf{Space}: $O(2^w)$ tile registers (classical full-state amplitude storage).
  \item \textbf{Time}: $O(d)$ global AFLG cycles, with constant-factor slack from routing latency absorbed into Wait$(\tau)$.
  \item \textbf{Error}: per-gate approximation error at most $\varepsilon$, yielding total variation distance from ideal output at most $O(d\varepsilon)$.
\end{itemize}
\end{proposition}

\begin{proof}[Proof sketch]
Space follows from Lemma~\ref{lem:embedding}: $2^w$ tiles suffice at depth $n = \lceil w / \log_2 k \rceil$. Time follows from Lemma~\ref{lem:gates}: each circuit layer maps to one global AFLG cycle, so $d$ layers cost $O(d)$ cycles. The constant-factor slack accounts for CDR phantom-to-physical link promotion and echo-aggregation overhead, both $O(n)$ substeps per cycle, giving total wall-clock $O(dn)= O(d \cdot w / \log_2 k)$. Error accumulation over $d$ sequential gates with independent per-gate error $\varepsilon$ is bounded by a union bound, giving total variation distance $\le d\varepsilon$.
\end{proof}

\subsection{Main Simulation Theorem}

\begin{theorem}[FTCA as Efficient Classical Quantum-Circuit Simulator]
\label{thm:main}
For any quantum circuit $C \in \mathcal{C}_{w,d}$ over the universal gate set $\{H, T, \mathrm{CNOT}\}$ and any $\varepsilon > 0$, there exists a fractal depth $n = \Theta(w / \log_2 k)$ and an FTCA program at depth $n$ such that a sequence of $O(d)$ global AFLG cycles computes a classical state-vector evolution whose final amplitude distribution approximates that of $C$ up to total variation distance $\le \varepsilon$. The space overhead is $O(2^w)$ tile registers, and the per-gate approximation error is at most $\varepsilon / d$.
\end{theorem}

\begin{proof}
Combine Lemma~\ref{lem:embedding} (state encoding), Lemma~\ref{lem:gates} (gate gadgets), and Proposition~\ref{prop:scaling} (resource bounds), setting $\varepsilon_{\text{gate}} = \varepsilon / d$, achievable with $p = \lceil \log_2(d/\varepsilon) \rceil$ bits of DAC precision.
\end{proof}

\begin{remark}
Theorem~\ref{thm:main} is a statement about \emph{classical} universality: FTCA performs full-state amplitude simulation in parallel, with the exponential state space stored explicitly across $O(2^w)$ tiles. It does not assert any speedup over conventional classical simulation; rather, it establishes that the FTCA substrate is a valid and efficient implementation model for quantum circuit simulation. The FSEP ontological interpretation is treated separately below.
\end{remark}

\subsection{Emergent Quantum Mechanics (FSEP Corollary)}

The following corollary separates the ontological claim from the simulation theorem above. It is conditional on the FSEP ontology and is not needed for Theorem~\ref{thm:main}.

\begin{corollary}[Emergent QM Realization]
\label{cor:fsep}
Under the FSEP ontology, where quantum mechanics itself arises as boundary resonances of the Apollonian substrate, the internal optical dynamics of FTCA at each AFLG interaction depth are isomorphic to those same M\"obius-induced resonances. Hence the same FTCA configuration both classically simulates a given quantum circuit (Theorem~\ref{thm:main}) and physically realizes an emergent-quantum evolution satisfying the same algebraic structure, up to hardware precision.

Formally: the light-field dynamics of FTCA are governed by the same Apollonian group action $\mathcal{A} = \langle S_1, S_2, S_3, S_4 \rangle$ and M\"obius algebra as the FSEP substrate. If FSEP is correct, then every quantum computation performed on FTCA is a faithful physical realization of the algebraic processes that generate quantum behavior in that ontology. The machine is thus not only a simulator---it is a \emph{self-simulating} quantum substrate.
\end{corollary}

\begin{proof}[Proof sketch]
The FTCA hardware geometry ($\lambda_{\text{local}} \approx 21.81$ from pixel aperture design) is chosen to match the FSEP Apollonian packing parameter. Each AFLG optical interaction corresponds to a discrete M\"obius transformation in $\mathcal{A}$; the echo-shell aggregation on QUANTUM\_STATE channels implements the boundary resonance summation that FSEP identifies with wavefunction collapse. Born-rule statistics are preserved to machine precision because amplitude norms are maintained by the echo-aggregation normalization protocol of Lemma~\ref{lem:embedding}. The parameter-free claim---that $\alpha^{-1}$ and other constants are determined by the fractal geometry---belongs to FSEP and is established in the cosmological derivation sections; here we only note that FTCA inherits those constants by construction.
\end{proof}

\subsection{Shared Algebraic Substrate}

Both FSEP and FTCA are governed by the \textbf{Apollonian group} $\mathcal{A} = \langle S_1, S_2, S_3, S_4 \rangle$, where each generator $S_i$ is a spherical inversion (a M\"obius transformation in $\mathbb{R}^3 \cup \{\infty\}$):
\[
S_i(\mathbf{r}) = \mathbf{c}_i + \frac{R_i^2 (\mathbf{r} - \mathbf{c}_i)}{\|\mathbf{r} - \mathbf{c}_i\|^2}.
\]
In \textbf{FSEP}, these generate the persistent fractal boundaries; residual structure at recursion depth $m$ produces effective quantization via boundary resonances. In \textbf{FTCA}, the same generators label the AGA $(k, w)$, with $w = \prod p_{i_j}$ encoding the reduced word in $\mathcal{A}$; prime factorization recovers the exact group element; word length equals fractal depth.

The Hilbert space basis for $w$ qubits is canonically embedded by interpreting generator words of length $\le w$ as binary strings (via prime-to-bit mapping modulo symmetry permutations). Thus the state vector $|\psi\rangle = \sum_{i} \alpha_i |i\rangle$ lives on the fractal skeleton whose connectivity graph $G_n$ (vertices = T2 Balls, edges = tangency + tube links) has adjacency spectrum approximating the Laplacian on the Apollonian limit set (Hausdorff dimension $D \approx 2.473946$).

\subsection{Fractal Replication as Multi-Scale Quantum Evolution}

The AFLG primitive replicates identically at every scale (pixel $\to$ tile $\to$ T2 Ball $\to$ VOC Ball $\to$ VOS):
\[
\text{AFLG}^{(m+1)} = \text{aggregate}\bigl(\{\text{AFLG}^{(m)}_j\}_{j \in \text{subspheres}}\bigr),
\]
where aggregation uses phantom CDR entries at VOC/VOS tangencies. This mirrors the renormalization-group flow induced by the FSEP M\"obius + scale-flip map $r \mapsto r/\lambda$.

In quantum information terms, this realizes a fractal tensor-network contraction:
\[
|\psi_{m+1}\rangle = \mathrm{Tr}_{\partial B_m} \Bigl( \bigotimes_{j \in \text{children}} U_j^{(m)} |\psi_m\rangle \Bigr),
\]
with boundary partial trace performed optically via echo-shell summation on QUANTUM\_STATE channels (deterministic T3 fiber timing guarantees causality).

\paragraph{Summary.} This section has established three independent results: Theorem~\ref{thm:main} gives a referee-checkable classical simulation theorem; Proposition~\ref{prop:scaling} quantifies the resource overhead; and Corollary~\ref{cor:fsep} connects the simulator to the FSEP ontology as a conditional emergent-QM claim. The construction satisfies the MATHICCS persistence requirement (eternal AFLG cycles at every recursion depth), and universal quantum simulation emerges as a direct consequence of the underlying M\"obius geometry. Open challenge: any other MATHICCS-valid ontology must likewise embed universal quantum simulation, or fail the self-consistency criterion.

\end{document}